\scanpage{80}
\[ s_0,\ s_1=\sigma_0s_0\ (=us_0=\sigma_0\underset{s_0}{\underbrace{\sigma_1s_0}}),\ s_{-1}^{\unsure{1|\prime}}=\sigma_1s_1\ (=s_{-1}=\sigma_1\sigma_0s_0) \]
On a le

\uline{Lemme} Étant donnés (2) \deleted{\unsure{(2)}}(3) satisfaisant \deleted{\unsure{(7)}}(8),
les conditions suivantes sont équivalentes :

\begin{tabular}{|p{.91\linewidth}}(a) Ces données correspondent à une réalisation
géom de $\Pi_\infty$, i.e les $s_n=u^n s_0=(\sigma_0\unsure{r_1|\sigma_1})^n s_0$
sont affinement indépendants

(b) $s_0,s_1$ et $s_{-1}$ sont affinement indépendts.

(c) les applications $\varphi_s$ et $\varphi_a$ de (3) sont
l'une et l'autre non constantes.\end{tabular}

En effet, $b\Rightarrow(a)\Rightarrow(c)$ sont \unsure{trivials|triviales|triviaux}, prouvons
$c)\Rightarrow a)$. On suppose c), il faut prouver

\uline{1}°) $s_1\ne s_0$ sinon $s_0$ serait fixe par $\sigma_1$ \uline{et} $\sigma_0$,
donc par $G$, donc par les $u^n$, donc $\varphi_s$ serait
constante.

2°) $s_{-1}\notin\operatorname{dr}(s_0,s_1)$. En effet, \deleted{\unsure{comme}}
\deleted{\unsure{a comme|Comme} $\sigma_0$ \unsure{fi}}
\[ a_0=\operatorname{dr}(s_0,s_1) \tag*{\deleted{\unsure{(8)}}(9)} \]
car $s_0\in a_0$, et comme \deleted{\unsure{par}} \inserted{$\sigma_{\unsure{0|1}}a_0=a_0$}, on doit
avoir $\sigma_0s_0=s_1\in a_0$. \deleted{Comme} On a
\[ \sigma_0a_0=\operatorname{dr}(\sigma_0s_0\underset{,}{=}s_1\ \sigma_1s_1=s_{-1}) \]
et si on avait $s_{-1}\in a_1$, on aurait $\sigma_{\unsure{1|0}}a_0=a_0$,
% archive sheet number: 79
